% Generated draft: requires mathematical and source review.
\documentclass[11pt]{article}
\usepackage[margin=25mm]{geometry}
\usepackage{fontspec}
\setmainfont{Latin Modern Roman}
\usepackage{amsmath,amssymb,mathtools,unicode-math}
\setmathfont{Latin Modern Math}
\usepackage{enumitem,needspace,xurl,hyperref}
\setlength{\parindent}{0pt}
\setlength{\parskip}{5pt}
\setlength{\emergencystretch}{3em}
\newcommand{\N}{\mathbb{N}}
\newcommand{\Z}{\mathbb{Z}}
\newcommand{\Q}{\mathbb{Q}}
\newcommand{\R}{\mathbb{R}}
\newcommand{\C}{\mathbb{C}}
\newcommand{\F}{\mathbb{F}}
\newcommand{\A}{\mathbb{A}}
\newcommand{\PP}{\mathbb{P}}
\newcommand{\E}{\mathbb{E}}
\newcommand{\eps}{\varepsilon}
\newcommand{\dd}{\,\mathrm{d}}
\begin{document}
\section*{UnsolvedMath ID: 30003125}
\textbf{Source problem code:} catalogue entry
\subsection{Definitions and mathematical statement}
The supplied record calls this the mixed Littlewood conjecture for arbitrary pseudo-absolute values. Its formula requires conventions that are not explicit in the record: the entries of the sequence must have a specified domain, the range of the limiting variable must be specified, and the defining minimum needs an initial product. The following precise formulation uses the intended-looking integer interpretation; these conventions require mathematical review against the cited source.

Let \(D=(d_i)_{i\geq 1}\) be any sequence of integers with \(d_i\geq 2\). Set
\[
D_0=1,\qquad D_n=\prod_{i=1}^{n}d_i\quad(n\geq 1).
\]
For each nonzero integer \(a\), define
\[
|a|_D=\min\{D_n^{-1}:n\in\Z_{\geq 0},\ a\in D_n\Z\}.
\]
This minimum exists because the set of eligible indices contains \(0\) and is finite. For \(t\in\R\), write
\[
\|t\|=\min_{m\in\Z}|t-m|
\]
for the distance to the nearest integer.

Under these explicitly stated conventions, the question is whether, for every such sequence \(D\) and every \(x\in\R\),
\[
\liminf_{\substack{q\to\infty\\q\in\Z_{\geq 1}}}q\,|q|_D\,\|qx\|=0.
\]
Equivalently, is it true that for every such \(D\), every \(x\in\R\), every \(\varepsilon>0\), and every integer \(Q\geq 1\), there exists an integer \(q\geq Q\) satisfying
\[
q\,|q|_D\,\|qx\|<\varepsilon?
\]
No boundedness assumption on the sequence \((d_i)\) is imposed. The supplied definition does not give a minimum at \(a=0\); no value there is needed for this formulation. The integer domain, the convention \(D_0=1\), and the positive-integer range of \(q\) are explicit interpretive additions, not verified quotations from the source.
\subsection{Short English statement}
Under the explicitly stated integer and indexing conventions, the conjecture asks whether every real number admits arbitrarily large positive integers \(q\) for which the product of \(q\), its \(D\)-pseudo-absolute value, and the distance of \(qx\) to the nearest integer is arbitrarily small. The source labels the problem open; that label and the formulation require independent review.
\subsection{Sources}
\begin{itemize}
\item Local UnsolvedMath registry: \texttt{data/unsolvedmath/problems.json}, permanent ID \(30003125\), source problem code \texttt{catalogue entry}, title \emph{The Mixed Littlewood Conjecture for Arbitrary Pseudo-Absolute Values}. The supplied record gives proposed year 2016 and does not name a proposer.
\item Supplied report citation: \emph{Diophantische Approximationen (2017)}, Oberwolfach Reports, pp.\ 1099--1187, DOI \texttt{10.4171/owr/2016/21}: \url{https://doi.org/10.4171/owr/2016/21}. The background describes these reports as records of mathematical workshops and their open-problem discussions.
\item Supplied literature source: Badziahin, \emph{Oberwolfach Report on Mixed Littlewood}: \url{https://doi.org/10.4171/owr/2016/21}. This attribution is supplied for the report reference and does not identify a proposer of the conjecture.
\item Supplied literature source: Badziahin and Harrap, \emph{Cantor-Winning Sets and Their Applications}: \url{https://arxiv.org/abs/1503.04738}.
\item Literature assessment dated exactly \texttt{2026-08-21}: the supplied curation labels the conjecture open in full generality and states that existing results provide special cases and restrictions, without a proof or counterexample for all sequences \(D\). This is supplied dated curation derived from the Oberwolfach Reports collection, not an independently verified literature assessment. Its time-sensitive status and cited claims require independent verification.
\end{itemize}
\end{document}
